\section{四个对照模型及比较原则}
\subsection{Equal Weight}
Equal Weight 在每个月度决策日，将资本平均分配给当期合格资产。若合格资产为 $n_t$ 个，则目标为 $w_i=1/n_t$。这种方法不需要收益预测，也不需要将协方差矩阵输入分配公式。资产的上市及准入仍由历史窗口决定，因此等权对象随历史可用性变化。

等权的优点在于规则透明，结果差异容易与资产池结构联系。在本研究中，多只权益和商品工具并存，等资本会使高波动资产承担较多组合风险。等权收益较高时，不能据此推断其风险控制更好；同样，较大回撤也不意味着等权失去参照价值。它直接呈现持有合格资产集合本身所带来的资本暴露。

\subsection{60/40 Benchmark}
60/40 Benchmark 按代码的资产分类，将 60\% 资本分给权益组，40\% 分给债券组，并在各组内等权。设合格权益与债券数量分别为 $n_E,n_B$，则
\[
w_i=\begin{cases}
0.60/n_E,&i\in E,\\
0.40/n_B,&i\in B,\\
0,&\text{其他资产}.
\end{cases}
\]
该对照是固定资本预算，不是风险预算。两组的实际风险贡献会随着波动与相关性变化，权益风险可能明显超过其资本比例。若某一期缺少必要资产组，则程序依既有规则处理，而不使用未来资产补足组合。

\subsection{HRP Benchmark}
HRP Benchmark 先根据合格资产的样本协方差构造相关矩阵，再将相关性转换为距离
\[
d_{ij}=\sqrt{\frac{1-\rho_{ij}}2}.
\]
程序采用单链接层次聚类获取叶节点顺序，将排序后的资产递归二分。对于一个簇，以簇内方差倒数构造临时权重，再计算该簇方差。两个子簇按对方方差占方差总和的比例分配资本，较低方差的一侧得到较高份额。

该分配过程使用相关性来组织资产，同时使用簇方差确定资本比例。它不等价于全局求解同一风险预算目标，也没有直接表达收益短缺的变量。实际结果受距离、链接方式、排序与二分规则共同影响，不能只用“不需要矩阵逆”概括其全部行为。

低波动工具可能使层次配置呈现较高资本集中度。对其高夏普，需要同时观察收益的绝对水平和波动分母。与 Global RRP 比较时，评价的重点是完整配置规则形成的风险收益位置，而非只对夏普作单一排名。

\subsection{HERC Benchmark 的实现范围}
仓库的 HERC Benchmark 与 HRP Benchmark 共用相关性距离和资产排序。区别在于递归分配阶段使用子簇波动而非方差比例。若子簇方差为 $v_L,v_R$，左侧资本比例为
\[
\alpha_L=\frac{\sqrt{v_R}}{\sqrt{v_L}+\sqrt{v_R}},\qquad
\alpha_R=1-\alpha_L.
\]
这种规则相对于方差倒数分配减弱了对子簇波动差异的放大，构成明确可复现的层次对照。它是本仓库的简化递归实现，不包含原始文献中所有聚类选择和风险预算步骤。

\subsection{共同条件与解释边界}
五个模型使用相同评价日期、成本费率和零无风险利率口径。全部结果由保存的日收益统一汇总，避免分别调用不一致的年化公式。主模型周频而对照月频，这一设计满足当前研究定位，但不支持从结果中识别纯粹的频率因果效应。

本文不展示其他 Global RRP 配置的测试结果，也不借对照的某一项劣势证明指定规格全面最优。高收益、低波动、低回撤和低换手通常难以同时占优。读者应把表格视为同一资产环境中的不同实施方案，而非一项已经排除了所有替代方案的最优性结论。
